\documentclass[11pt]{article}
\usepackage[T1]{fontenc}
\usepackage[utf8]{inputenc}
\usepackage{amsmath,amssymb,amsthm}
\usepackage{booktabs}
\usepackage[margin=1in]{geometry}
\usepackage[hidelinks]{hyperref}

\newcommand{\Var}{\operatorname{Var}}
\newcommand{\R}{\mathbb{R}}
\newcommand{\EE}{\mathbb{E}}

\newtheorem{theorem}{Theorem}[section]
\newtheorem{lemma}[theorem]{Lemma}
\newtheorem{proposition}[theorem]{Proposition}
\newtheorem{corollary}[theorem]{Corollary}
\newtheorem{conjecture}[theorem]{Conjecture}
\theoremstyle{definition}
\newtheorem{definition}[theorem]{Definition}
\newtheorem{problem}[theorem]{Open problem}
\theoremstyle{remark}
\newtheorem{remark}[theorem]{Remark}

\title{Equichordal Points on Ellipses and on an Odd-Harmonic Family:\\
Proofs, Bounded Scans, and the Limits of Both}
\author{Research manuscript prepared within a multi-agent workflow\\
\small proofs, numerical readings and literature sources were produced by different agents of the same workflow}
\date{}

\begin{document}
\maketitle

\begin{abstract}
An equichordal point of a plane region is an interior point such that all chords through it have one and the same
length. We give a self-contained account of what is known, what can be proved elementarily, and what has only
been scanned, for two families: the ellipse $E(a,b)$ with $a\neq b$ and a radial family obtained by perturbing the
unit circle with two odd harmonics.

The mathematical part contains nine statements with complete proofs. Three are structural: the chord-length
functional is continuous in the point and in the direction, so its standard deviation has a modulus of
continuity; a similarity of the plane multiplies the variance of the chord length by the square of the scale,
so only the shape ratio matters; and a positive-definite Hessian at a critical point gives an explicit ball on
which the variance cannot vanish. Two are classification results: a disc has exactly one equichordal point,
its centre, and the centre of a non-circular ellipse is never one. One is a construction: if the radial function
satisfies $r(\theta)+r(\theta+\pi)\equiv 2c$, in particular for any finite sum of odd harmonics with
$\varepsilon<1$, the origin is an equichordal point with chord length exactly $2c$. Three concern the limits of
computation: an a priori error bound for the bisection evaluation of chord lengths, a covering-radius criterion
that a finite grid must satisfy before it can certify the absence of a zero of the variance, and the quantified
outcome of the finite scans themselves.

The numerical part reports scans of nine shape ratios in $[1.01,5]$ and four values of the perturbation
parameter, with two direction counts and two evaluation precisions; in none of them did a candidate away from
the known exact points meet the pre-registered threshold. Both family-wide statements that these scans suggest
are stated as conjectures, not theorems, and the discussion lists the seven specific gaps between the scans and
their conjectures. No machine-checked proof exists for anything in this manuscript, and no statement here
claims one.
\end{abstract}

\section{Introduction}

\subsection{The problem}

Let $\Gamma$ be a simple closed plane curve bounding a region $K$, and let $P$ be an interior point of $K$. The
point $P$ is called an \emph{equichordal point} of $\Gamma$ if there is a constant $L>0$ such that every line
through $P$ meets $\Gamma$ in two points whose distance is $L$; the curve is then an equichordal curve and $L$
is the chord length of $P$. The terminology is that of the problem collection of Klee \cite{klee1960}, whose
formulation we quote verbatim in Section~\ref{sec:intro-lit} because it is stated for star-shaped sets and not
only for convex bodies.

The equichordal point problem asks whether a plane convex body can possess two distinct equichordal points. The
question was asked by Fujiwara in 1916 and, independently, by Blaschke, Rothe and Weitzenb\"ock in 1917, as the
quoted commentary records. The answer is negative for plane convex bodies: the complete solution is due to
Rychlik \cite{rychlik1997}, who reduced the problem to the non-existence of heteroclinic connections for a
multi-valued algebraic map. A disc shows that one equichordal point is easy: its centre is one, with chord
length equal to the diameter. The whole difficulty lies in the second point.

\subsection{Known results and the exact class they cover}\label{sec:intro-lit}

We reproduce only what can be pointed at with a verbatim quotation and a page of the consulted reprint; printed
page numbers are given, and quotations are literal, including the source's spelling of names.

The definition, printed p.\,24:
\begin{quote}
``If the set $C$ is star-shaped from its interior point $p$ (i.e. $C$ contains every segment joining $p$ to
another point of $C$), $p$ is called an equichordal point of $C$ provided all chords of $C$ which pass through
$p$ are of the same length.'' \cite{klee1960}
\end{quote}
Two features of this wording matter below. The definition is given for star-shaped sets, so it is not a
convexity-only notion; and it is a property of a point with respect to a set, not of a curve alone.

The attribution, printed p.\,24:
\begin{quote}
``Whether a plane convex body can possess two equichordal points was first asked by Fujiwara [5], and
independently by Blaschke, Rothe, and Weitzenbock [1].'' \cite{klee1960}
\end{quote}
This fixes the origin of the question and asserts nothing about its answer.

The accepted reading of the 1997 solution, printed p.\,27, in a commentary attached to the same problem:
\begin{quote}
``The accepted wisdom is that the 1997 paper [Ryc97] of Rychlik finally disposed of the problem, by showing that
no 2e-curve (of a type that includes convex curves) can have two equichordal points.'' \cite{klee1960}
\end{quote}
Two restrictions are carried by this sentence and we keep both. First, it is a commentary's paraphrase, not a
quotation from \cite{rychlik1997} itself, and we use it only as a second-hand report of a convention. Second,
the class is described as ``a type that includes convex curves'', which is not the class of all simple closed
curves; a statement about arbitrary simple closed curves is therefore not covered. A further commentary
sentence, printed p.\,27, reads:
\begin{quote}
``In the early articles about the problem it is established that if a 2e-curve exists, it must have certain
symmetry and smoothness properties, regardless of whether it is assumed to be convex or not; that much is not
in dispute.'' \cite{klee1960}
\end{quote}
We take from it only that certain symmetry and smoothness properties are available without a convexity
hypothesis, so such properties must not be cited as convex-only.

\subsection{What this manuscript adds and what it does not}

The two families studied here are deliberately special, because they are the ones on which a positive answer can
be tested cheaply.
\begin{enumerate}
\item \emph{The ellipse.} For $E(a,b)$ with $a\neq b$: is any interior point equichordal? Nothing in the quoted
literature answers this, since the quoted negative result concerns \emph{two} equichordal points on a curve of a
restricted class, not one point on an ellipse. Section~\ref{sec:main} proves that the centre is not one
(Proposition~\ref{prop:ellipse-centre}), and Section~\ref{app:comp} reports the scans.
\item \emph{An odd-harmonic radial family.} Perturbing the circle by two odd harmonics produces a family whose
origin is an exact equichordal point, with chord length exactly $2$ and with a proof rather than a reading
(Theorem~\ref{thm:odd-harmonic}). The question is then whether a second one exists, and that question is answered
here only by bounded scans.
\end{enumerate}
We claim no result about general convex bodies, about general non-convex simple closed curves, about shape ratios
above $5$, about the limit ratio $1$, or about the limit of vanishing perturbation. Those exclusions are listed
with their reasons in Section~\ref{sec:discussion} and, where the obstruction is a resolution limit of the
method rather than a missing computation, they are stated as propositions about the method
(Proposition~\ref{prop:grid-certificate}). Every numerical statement in this manuscript is a statement about a
finite grid, a finite direction set and a stated threshold, and is written as such.

Nothing in this manuscript is machine-checked. The proofs below are ordinary natural-language proofs written to
be read line by line; the numerical readings are outputs of the algorithm described in
Section~\ref{app:comp}. No claim is made that any proof has been verified by a computer, and no claim is made
that a scan establishes non-existence.

\section{Definitions and notation}\label{sec:defs}

\begin{definition}[Equichordal point]\label{def:eqp}
Let $\Gamma$ be a simple closed curve bounding a region $K$ with non-empty interior, and let $P\in\operatorname{int}K$.
For a line $\ell$ through $P$ meeting $\Gamma$ in exactly two points $A,B$ (with $P$ between them) write
$L_\ell=|AB|$. Then $P$ is an \emph{equichordal point} of $\Gamma$ if there is $L>0$ with $L_\ell=L$ for every
line $\ell$ through $P$. The constant $L$ is the \emph{chord length} of $P$, and $\Gamma$ is then an
\emph{equichordal curve}.
\end{definition}

We consider only curves for which every line through the relevant interior point meets the curve in exactly two
points. This holds, for instance, for every convex body and for every star-shaped region with a radial function
that is continuous and positive, which are the two cases used below.

\begin{definition}[Radial function]\label{def:radial}
A region $K$ containing the origin is \emph{star-shaped with respect to} $0$ if every ray from $0$ meets $\Gamma$ in
exactly one point. It is then described by its radial function $r:\R\to(0,\infty)$, $2\pi$-periodic and continuous,
through $\Gamma=\{r(\theta)(\cos\theta,\sin\theta):\theta\in\R\}$, and the chord through the origin in direction
$\theta$ has length $r(\theta)+r(\theta+\pi)$.
\end{definition}

\begin{definition}[Chord-length function and its variance]\label{def:var}
For $P$ interior and a unit vector $u(\varphi)=(\cos\varphi,\sin\varphi)$, let $L_P(\varphi)$ be the length of the
chord through $P$ in direction $u(\varphi)$, and put
\[
m(P)=\frac{1}{2\pi}\int_0^{2\pi}L_P(\varphi)\,d\varphi,\qquad
F(P)=\Var_\varphi L_P(\varphi)=\frac{1}{2\pi}\int_0^{2\pi}\bigl(L_P(\varphi)-m(P)\bigr)^2d\varphi .
\]
By Definition~\ref{def:eqp}, $P$ is equichordal if and only if $F(P)=0$.
\end{definition}

For computations we replace the circle of directions by a finite set and the integral by a mean. For $N\ge 2$ put
$\Phi_N=\{2\pi k/N:k=0,\dots,N-1\}$ and
\[
m_N(P)=\frac1N\sum_{\varphi\in\Phi_N}L_P(\varphi),\qquad
V_N(P)=\frac1N\sum_{\varphi\in\Phi_N}\bigl(L_P(\varphi)-m_N(P)\bigr)^2,\qquad \sigma_N(P)=\sqrt{V_N(P)}.
\]

\begin{definition}[The two families]\label{def:families}
The ellipse with semi-axes $a,b>0$ is $E(a,b)=\{(x,y):x^2/a^2+y^2/b^2=1\}$, and its shape is described by the
ratio $a/b$. Its radial function is $r(\varphi)=ab\,(b^2\cos^2\varphi+a^2\sin^2\varphi)^{-1/2}$. The second family is
\[
r(\theta)=1+\varepsilon\bigl(0.6111\cos 3\theta+0.3889\cos 5\theta\bigr),\qquad 0<\varepsilon<1,
\]
with coefficients normalised so that $0.6111+0.3889=1$; it is a radial perturbation of the unit circle containing
only the odd harmonics $k=3,5$. We shall refer to it as the odd control family.
\end{definition}

\begin{definition}[Search criterion used below]\label{def:criterion}
For a search domain $D$ and a point $P\in D$, we call $P$ a \emph{candidate} if both $V_{2000}(P)<10^{-18}$ and
$V_{8000}(P)<10^{-18}$. In the odd-harmonic family the origin is excluded by construction, so candidates are
sought in an annulus $D=\{P: \rho_-\le|P|\le\rho_+\}$; for the ellipse the whole disc $|P|\le\kappa\min(a,b)$ is
searched. The threshold $10^{-18}$ and the two direction counts were fixed before the runs; they are not adjusted
per family.
\end{definition}

The threshold is far above double-precision noise on the quantities involved: a value of the variance of order
$10^{-27}$ is attained at an exact equichordal point (the centre of a circle, the origin of the odd-harmonic
family), so $10^{-18}$ is fourteen orders of magnitude above what an exact zero computes to in this
implementation. This is what makes the criterion a meaningful positive test; Section~\ref{app:comp} explains why
it is not a meaningful negative test.

\section{Main results}\label{sec:main}

All numerical evidence referred to anywhere in this paper is collected in Appendix~A; conclusions whose justification depends on that evidence are marked as such in the text, and none of them is stated as a theorem.

Throughout, $\Gamma$ is a simple closed curve bounding a region $K$, and $L_P$ is as in Definition~\ref{def:var}.

\begin{lemma}[Continuity of the chord-length function]\label{lem:continuity}
Let $K$ be compact, star-shaped with respect to the origin and bounded by a curve with continuous positive radial
function $r$. Let $K^\circ$ denote its interior and let $D\Subset K^\circ$ be compact. Then:
\begin{enumerate}
\item[(i)] for every $P\in K^\circ$ the map $\varphi\mapsto L_P(\varphi)$ is continuous and $2\pi$-periodic;
\item[(ii)] $P\mapsto L_P$ is continuous from $D$ to $C^0(\R/2\pi\R)$ equipped with the supremum norm, hence
$F$ and $\sigma=\sqrt{F}$ are continuous on $D$; consequently $\sigma$ has a modulus of continuity $\omega_D$ on
$D$, that is, $|\sigma(P)-\sigma(Q)|\le\omega_D(|P-Q|)$ for all $P,Q\in D$.
\end{enumerate}
\end{lemma}

\begin{proof}
(i) Fix $P\in K^\circ$ and a direction $u(\varphi)$. The line $\{P+t\,u(\varphi):t\in\R\}$ contains $P$, which lies
in the open set $K^\circ$, and it meets $\Gamma$ in exactly two points because $K$ is star-shaped with a positive
continuous radial function. Write the corresponding parameters as $t_-(\varphi)<0<t_+(\varphi)$, so that
$L_P(\varphi)=t_+(\varphi)-t_-(\varphi)$. The function $G(t,\varphi)=|P+t\,u(\varphi)|-r(\arg(P+t\,u(\varphi)))$ is
continuous, is negative at $t=0$, and tends to $+\infty$ as $t\to\pm\infty$; its positive zero is exactly
$t_+(\varphi)$ and its negative zero exactly $t_-(\varphi)$. The set of zeros depends continuously on $\varphi$ in
the Hausdorff sense because $G$ is continuous and the zeros are simple ($\Gamma$ is a curve and the line crosses
it), so $t_+$ and $t_-$ are continuous, and with them $L_P$. Periodicity is clear.

(ii) Let $P,Q\in D$. Fix $\varphi$. The two chords are the intersections of two parallel lines with $\Gamma$. Both
lines pass through the compact set $D$, whose distance to $\Gamma$ is a positive number $\delta$. Moving the line
from $P$ to $Q$ moves each endpoint continuously along $\Gamma$, and $\Gamma$ is compact, so the map
$(P,\varphi)\mapsto L_P(\varphi)$ is uniformly continuous on $D\times\R/2\pi\R$. Hence for every $\eta>0$ there is
$\zeta>0$ with $|L_P(\varphi)-L_Q(\varphi)|\le\eta$ for all $\varphi$ whenever $|P-Q|\le\zeta$; that is,
$\|L_P-L_Q\|_\infty\le\eta$. Since the mean and the variance are continuous in the supremum norm (the mean is
$1$-Lipschitz and the standard deviation is also $1$-Lipschitz, as $|\sigma(X)-\sigma(Y)|\le\|X-Y\|_\infty$ for
random variables on a common probability space), $\sigma$ is continuous on the compact set $D$, hence uniformly
continuous, hence has a modulus of continuity.
\end{proof}

\begin{lemma}[Variance vanishes only for constant chord length]\label{lem:zero-var}
Let $P\in K^\circ$ and let $\varphi\mapsto L_P(\varphi)$ be continuous. Then $F(P)=0$ if and only if $L_P$ is
constant; that is, $P$ is equichordal.
\end{lemma}

\begin{proof}
\emph{Step 1 (integral of a continuous non-negative function).} The integrand $\varphi\mapsto(L_P(\varphi)-m(P))^2$ is
continuous and non-negative on the circle, so its integral vanishes if and only if it vanishes identically: a
continuous non-negative function that is positive at one point is positive on a neighbourhood of it, and that
neighbourhood has positive measure.

\emph{Step 2 (constant chord length).} Vanishing of the integrand is equivalent to $L_P(\varphi)=m(P)$ for all
$\varphi$, i.e.\ to $L_P$ being constant.

\emph{Step 3 (identification).} By Definition~\ref{def:eqp}, $L_P$ constant means that $P$ is an equichordal point.
\end{proof}

\begin{proposition}[The disc has exactly one equichordal point]\label{prop:disc}
Let $K$ be the closed disc of radius $R$ about $O$, and let $P$ be an interior point at distance $d$ from $O$.
Then, with $\psi$ the angle between the direction of the chord and the direction $OP$,
\[
L_P(\psi)=2\sqrt{R^2-d^2\sin^2\psi},\qquad
F(P)=4R^2\Bigl[\Bigl(1-\tfrac{k^2}{2}\Bigr)-\Bigl(\frac{2E(k)}{\pi}\Bigr)^{2}\Bigr],\qquad k=\frac{d}{R},
\]
where $E(k)=\int_0^{\pi/2}\sqrt{1-k^2\sin^2\psi}\,d\psi$ is the complete elliptic integral of the second kind.
Consequently $F(P)=0$ if and only if $d=0$: the centre is the only equichordal point of the disc.
\end{proposition}

\begin{proof}
\emph{Step 1: chord length.} A point of the line through $P$ in direction $u$ has the form $P+tu$, and it lies on
the circle when $|P+tu|^2=R^2$, that is
\[
t^2+2t\,(P\cdot u)+d^2-R^2=0 .
\]
The two roots $t_1<t_2$ satisfy $t_1+t_2=-2P\cdot u$ and $t_1t_2=d^2-R^2$, hence
\[
L_P=t_2-t_1=\sqrt{(t_1+t_2)^2-4t_1t_2}=\sqrt{4(P\cdot u)^2+4(R^2-d^2)}=2\sqrt{R^2-d^2+(P\cdot u)^2} .
\]
With $\psi$ the angle between $u$ and $OP$ we have $P\cdot u=d\cos\psi$, so
$L_P=2\sqrt{R^2-d^2+d^2\cos^2\psi}=2\sqrt{R^2-d^2\sin^2\psi}$, which is the first formula.

\emph{Step 2: two directions.} Taking $\psi=0$ gives $L_P=2R$, and taking $\psi=\pi/2$ gives
$L_P=2\sqrt{R^2-d^2}$. If $d>0$ these two values differ, so $L_P$ is not constant and
Lemma~\ref{lem:zero-var} gives $F(P)>0$.

\emph{Step 3: the centre.} If $d=0$ the formula of Step 1 gives $L_P\equiv 2R$, which is constant, so $F(P)=0$ and
$P$ is equichordal.

\emph{Step 4: exact variance.} Averaging the square of the Step~1 formula over $\psi$ uniform on $[0,2\pi)$ and using
$\int_0^{2\pi}\sin^2\psi\,d\psi=\pi$ gives $\EE[L_P^2]=4R^2(1-k^2/2)$. For the mean, substituting the four
quadrants into $E(k)$ gives $\EE[L_P]=(1/2\pi)\int_0^{2\pi}2R\sqrt{1-k^2\sin^2\psi}\,d\psi=\frac{4R}{\pi}E(k)$.
Hence $F(P)=\EE[L_P^2]-\EE[L_P]^2=4R^2[(1-k^2/2)-(2E(k)/\pi)^2]$. This formula is not needed for the
classification, which Steps 2 and 3 already give, but it shows how the variance of a disc grows from zero: at
$k=0$ the two terms coincide, and the expansion of $E$ shows that the bracket is of order $k^4$.
\end{proof}

\begin{proposition}[The centre of a non-circular ellipse is not an equichordal point]\label{prop:ellipse-centre}
Let $a>b>0$ and let $E(a,b)$ be the ellipse with semi-axes $a,b$. Then its centre $O$ is not an equichordal point;
more precisely $F(O)>0$.
\end{proposition}

\begin{proof}
\emph{Step 1: radial function.} A point of $E(a,b)$ in direction $\varphi$ has the form $\rho(\cos\varphi,\sin\varphi)$
with $\rho^2(\cos^2\varphi/a^2+\sin^2\varphi/b^2)=1$, hence
\[
\rho(\varphi)=\frac{ab}{\sqrt{b^2\cos^2\varphi+a^2\sin^2\varphi}} .
\]
This is continuous and positive, and the ellipse is convex, so every line through an interior point meets it in
exactly two points.

\emph{Step 2: chords through the centre.} The chord through $O$ in direction $\varphi$ has endpoints
$\rho(\varphi)(\cos\varphi,\sin\varphi)$ and $-\rho(\varphi+\pi)(\cos\varphi,\sin\varphi)$, and
$\rho(\varphi+\pi)=\rho(\varphi)$; the chord length is therefore $L_O(\varphi)=2\rho(\varphi)$.

\emph{Step 3: two directions.} At $\varphi=0$ we get $L_O(0)=2\rho(0)=2a$; at $\varphi=\pi/2$ we get
$L_O(\pi/2)=2b$. Since $a\neq b$ the function $L_O$ is not constant. It is continuous by Step 1, so
Lemma~\ref{lem:zero-var} gives $F(O)>0$, and by Definition~\ref{def:eqp} the centre is not equichordal.
\end{proof}

\begin{remark}
Proposition~\ref{prop:ellipse-centre} says only that the obvious candidate fails. It does not say that no other
point of the ellipse is equichordal; that stronger statement is Conjecture~\ref{conjecture-ellipse} below, and the
scans of Section~\ref{app:comp} are what supports it.
\end{remark}

\begin{lemma}[Similarity and reflection invariance]\label{lem:similarity}
Let $a,b>0$, $\lambda>0$, and let $F_{a,b}$ denote the variance functional of Definition~\ref{def:var} for the
ellipse $E(a,b)$. Then
\[
F_{\lambda a,\lambda b}(\lambda P)=\lambda^2 F_{a,b}(P)\quad\text{for all }P\in\operatorname{int}E(a,b),
\]
and $F_{a,b}$ is invariant under the reflections $(x,y)\mapsto(-x,y)$ and $(x,y)\mapsto(x,-y)$. Consequently, the
sign of $F_{a,b}$ and the set of equichordal points depend on the ratio $a/b$ only, and the search for them may be
restricted to a closed quadrant of the plane.
\end{lemma}

\begin{proof}
\emph{Step 1: scaling.} The map $X\mapsto\lambda X$ sends $E(a,b)$ onto $E(\lambda a,\lambda b)$ and sends lines to
lines, so the chord through $\lambda P$ in direction $\varphi$ is the image of the chord through $P$ in direction
$\varphi$; hence $L^{\lambda a,\lambda b}_{\lambda P}(\varphi)=\lambda L^{a,b}_P(\varphi)$ for all $\varphi$.
Therefore $m(\lambda P)=\lambda m(P)$ and
\[
F_{\lambda a,\lambda b}(\lambda P)=\frac{1}{2\pi}\int_0^{2\pi}\bigl(\lambda L^{a,b}_P(\varphi)-\lambda m(P)\bigr)^2d\varphi
=\lambda^2F_{a,b}(P).
\]
\emph{Step 2: reflections.} The reflection $\varrho_1(x,y)=(-x,y)$ maps $E(a,b)$ onto itself and maps the chord
through $P$ in direction $(\cos\varphi,\sin\varphi)$ onto the chord through $\varrho_1P$ in direction
$(-\cos\varphi,\sin\varphi)$; the two chords have the same length, and the map on the direction circle is a
measure-preserving bijection, so the mean and the variance are unchanged. The same argument applies to
$\varrho_2(x,y)=(x,-y)$.

\emph{Step 3: consequences.} By Step 1 the equation $F_{\lambda a,\lambda b}(Q)=0$ is equivalent to
$F_{a,b}(Q/\lambda)=0$, so the answer does not change along a ray in the parameter $(a,b)$; it depends on $a/b$
only. By Step 2 the set $\{P:F_{a,b}(P)=0\}$ is invariant under both reflections, so it is determined by its
intersection with any closed quadrant.
\end{proof}

\begin{theorem}[Odd-harmonic radial families have the origin as an equichordal point]\label{thm:odd-harmonic}
Let $\varepsilon>0$, let $a_3,\dots,a_{2m+1}\in\R$ be finitely many coefficients with
$\sum_{j}|a_{2j+1}|=1$, and let
\[
r(\theta)=1+\varepsilon\sum_{j=0}^{m}a_{2j+1}\cos\bigl((2j+1)\theta\bigr).
\]
Then:
\begin{enumerate}
\item[(i)] if $0<\varepsilon<1$ then $r(\theta)>0$ for all $\theta$, so $r$ defines a star-shaped region with a
simple closed boundary curve;
\item[(ii)] $r(\theta)+r(\theta+\pi)=2$ for all $\theta$;
\item[(iii)] every chord through the origin has length exactly $2$, so the origin is an equichordal point with
chord length $2$.
\end{enumerate}
In particular this holds for $r(\theta)=1+\varepsilon(0.6111\cos3\theta+0.3889\cos5\theta)$ with $0<\varepsilon<1$, the
odd control family of Definition~\ref{def:families}.
\end{theorem}

\begin{proof}
\emph{Step 1: positivity.} For all $\theta$, $|r(\theta)-1|\le\varepsilon\sum_j|a_{2j+1}|=\varepsilon$, so
$r(\theta)\ge1-\varepsilon>0$. A continuous positive $2\pi$-periodic radial function defines a simple closed curve
that is star-shaped with respect to the origin, and every ray from the origin meets it once.

\emph{Step 2: the antipodal relation.} For every integer $k$, $\cos(k(\theta+\pi))=(-1)^k\cos(k\theta)$. All the
frequencies in the sum are odd, so each term changes sign, and
\[
r(\theta+\pi)=1-\varepsilon\sum_j a_{2j+1}\cos\bigl((2j+1)\theta\bigr)=2-r(\theta).
\]

\emph{Step 3: chord lengths through the origin.} Fix $\theta$ and let $u=(\cos\theta,\sin\theta)$. The ray from the
origin in direction $u$ meets the curve at the point $r(\theta)u$, and the opposite ray meets it at the point
$-r(\theta+\pi)u$; both distances are positive by Step 1, and the two points are on the same line through the
origin, on opposite sides of it. Hence the chord through the origin in the direction $u$ has length
\[
r(\theta)+r(\theta+\pi)=r(\theta)+2-r(\theta)=2 ,
\]
independently of $\theta$. By Definition~\ref{def:eqp} the origin is an equichordal point with chord length $2$, and
by Lemma~\ref{lem:zero-var} $F(0)=0$.
\end{proof}

\begin{remark}
Theorem~\ref{thm:odd-harmonic} shows that the property ``has an equichordal point'' is cheap: it is implied by one
linear relation among the odd Fourier coefficients, namely $r(\theta+\pi)=2-r(\theta)$. Nothing similar is known
to us for a second equichordal point, and the theorem says nothing in that direction. Note also that the family is
not contained in the class of curves for which the quoted negative result is stated, so no contradiction with
Section~\ref{sec:intro-lit} is possible: the theorem produces one equichordal point, the quoted result excludes
two.
\end{remark}

Next we record exactly what a finite grid can and cannot certify. The statement is about the method; it is the
mathematical reason why the scans of Section~\ref{app:comp} support a conjecture rather than a theorem.

\begin{proposition}[Grid certificate and its covering-radius condition]\label{prop:grid-certificate}
Let $D\Subset\operatorname{int}K$ be compact, let $\sigma=\sqrt{F}$ on $D$ with modulus of continuity
$\omega_D$ as in Lemma~\ref{lem:continuity}, and let $G\subset D$ be a finite set with covering radius
$\rho=\sup_{P\in D}\min_{g\in G}|P-g|$. If
\[
\min_{g\in G}\sigma(g)>\omega_D(\rho),
\]
then $F(P)>0$ for every $P\in D$; that is, $D$ contains no equichordal point. Conversely, if $D$ contains an
equichordal point then necessarily $\min_{g\in G}\sigma(g)\le\omega_D(\rho)$.
\end{proposition}

\begin{proof}
\emph{Step 1.} Let $P\in D$ be arbitrary and let $g\in G$ satisfy $|P-g|\le\rho$; such $g$ exists by the definition
of covering radius. By Lemma~\ref{lem:continuity}, $\sigma(P)\ge\sigma(g)-\omega_D(|P-g|)\ge\sigma(g)-\omega_D(\rho)$,
using that a modulus of continuity is non-decreasing.

\emph{Step 2.} If $\sigma(g)>\omega_D(\rho)$ for every $g\in G$, Step 1 gives $\sigma(P)>0$ for every $P\in D$, hence
$F(P)>0$ on $D$ and, by Lemma~\ref{lem:zero-var}, no point of $D$ is equichordal.

\emph{Step 3.} For the converse, if $P_0\in D$ were equichordal then $\sigma(P_0)=0$, and Step 1 applied to
$P_0$ gives $0=\sigma(P_0)\ge\sigma(g_0)-\omega_D(\rho)$ for a nearest grid point $g_0$, that is
$\sigma(g_0)\le\omega_D(\rho)$, so the minimum over $G$ is at most $\omega_D(\rho)$.
\end{proof}

The content of Proposition~\ref{prop:grid-certificate} is that a scan certifies absence only by the margin between
the smallest observed standard deviation and the modulus of continuity at the covering radius. A coarse grid has
a large covering radius and therefore a large $\omega_D(\rho)$; a fine local refinement lowers the observed
minimum but does not by itself lower $\omega_D(\rho)$ for the region between the grid points. Both effects are
visible in the numbers of Section~\ref{app:comp}, where the smallest observed $\sigma$ is of order $3.5\times10^{-5}$
while $\omega_D(\rho)$ for any grid used there is of order the grid spacing itself.

\begin{lemma}[A positive-definite Hessian gives an explicit ball of non-vanishing]\label{lem:hessian}
Let $F$ be twice continuously differentiable on a ball $B(P_0,R_0)$ with $\nabla F(P_0)=0$ and
$D^2F(P_0)\succeq \lambda I$ for some $\lambda>0$. Suppose in addition that $\|D^3F\|\le M$ on $B(P_0,R_0)$ in the
sense that $|F(P)-F(P_0)-\tfrac12\langle D^2F(P_0)(P-P_0),P-P_0\rangle|\le \tfrac{M}{6}|P-P_0|^3$. Then for
\[
|P-P_0|\le \min\Bigl\{R_0,\ \frac{3\lambda}{2M}\Bigr\}
\]
one has $F(P)\ge F(P_0)+\tfrac{\lambda}{4}|P-P_0|^2$. In particular $F$ has no zero in that ball whenever
$F(P_0)>0$.
\end{lemma}

\begin{proof}
\emph{Step 1: Taylor.} By the second-order Taylor theorem with the stated bound on the third derivative,
\[
F(P)=F(P_0)+\tfrac12\langle D^2F(P_0)(P-P_0),P-P_0\rangle+R_3(P),\qquad |R_3(P)|\le\tfrac{M}{6}|P-P_0|^3 .
\]
\emph{Step 2: lower bound.} Write $h=P-P_0$ and $r=|h|$. By positive definiteness,
$\langle D^2F(P_0)h,h\rangle\ge\lambda r^2$, hence
\[
F(P)\ge F(P_0)+\tfrac{\lambda}{2}r^2-\tfrac{M}{6}r^3=F(P_0)+r^2\Bigl(\tfrac{\lambda}{2}-\tfrac{M}{6}r\Bigr).
\]
\emph{Step 3: the radius.} If $r\le 3\lambda/(2M)$ then $\tfrac{M}{6}r\le\tfrac{\lambda}{4}$, so
$\tfrac{\lambda}{2}-\tfrac{M}{6}r\ge\tfrac{\lambda}{4}$ and $F(P)\ge F(P_0)+\tfrac{\lambda}{4}r^2$. The last claim
follows since $F(P_0)>0$.
\end{proof}

Lemma~\ref{lem:hessian} is exactly what a recorded positive-definite Hessian at a numerically located minimiser
buys: a zero-free ball around that minimiser, not a zero-free domain. It is used in
Section~\ref{app:comp} to say precisely what the local rechecks establish and what they do not.

\long\def\CompEvidenceBody{%
% Body of Appendix~A; typeset at the end of the document, after the References.
% The brace that closes this definition stands immediately before Section~4.

\subsection{The algorithm as a mathematical specification}

The objects evaluated are the quantities of Definition~\ref{def:var} and the criterion of
Definition~\ref{def:criterion}. For an ellipse we set $b=1$ and vary $a$, which is legitimate by
Lemma~\ref{lem:similarity}; the search domain is the disc $D_\kappa=\{P:|P|\le\kappa\min(a,b)\}$ with
$\kappa\in\{0.6,0.9,0.99,0.999\}$ in the various runs, and the full interior $E(a,b)^\circ$ in one run.

The procedure at each stage is the following.
\begin{enumerate}
\item[(A1)] \emph{Global stage.} An $11\times11$ (or $13\times13$) uniform grid on the square circumscribing
$D_\kappa$, truncated to $D_\kappa$, together with the origin, provides the starting points. The covering radius of
this grid inside the square is $\sqrt{2}\,h/2$ with $h$ the grid spacing.
\item[(A2)] \emph{Local stage.} From each start a compass search minimises $F$, i.e.\ it evaluates the four (or
eight) neighbouring points on a square lattice of step $s$, moves to the best one if it is strictly better, and
divides $s$ by two when no neighbour improves, down to the step floor $s_{\min}=10^{-8}$ (ellipses) or $10^{-9}$
(odd-harmonic family).
\item[(A3)] \emph{Evaluation.} The chord lengths entering $V_N$ are computed by bisection ray shooting with $45$
iterations, and the candidate minimisers are re-evaluated with $60$ iterations; at the argmin of several runs a
fiftieth-digit evaluation by an independent multiple-precision implementation is used as a check on the
double-precision value.
\item[(A4)] \emph{Verdict.} A point is a candidate in the sense of Definition~\ref{def:criterion} if both
$V_{2000}$ and $V_{8000}$ fall below $10^{-18}$ there; the run reports the smallest value found, its argument, and
whether the argmin is interior or on the boundary of the search domain.
\end{enumerate}

Two quantities decide how a reading may be read. The first is the argument being interior or on the boundary: a
boundary argmin is a statement about the disc $D_\kappa$ only. The second is the covering radius of the global
stage, which by Proposition~\ref{prop:grid-certificate} is the only thing that could turn a scan into a
certificate of absence.

\begin{lemma}[A priori error bound for the bisection evaluation]\label{lem:bisection}
Let the chord lie in a disc of radius $R_{\max}$ about the origin and let each of its two endpoints be computed by
bisection to $n$ halvings, starting from an interval of length at most $2R_{\max}$. Then the computed chord
length differs from the true one by at most $4R_{\max}2^{-n}$.
\end{lemma}

\begin{proof}
\emph{Step 1: bracket length.} Starting from an interval of length $2R_{\max}$ containing the endpoint, after $n$
halvings the surviving interval has length $2R_{\max}2^{-n}$, and the reported endpoint lies in it, so each
endpoint is known to within $2R_{\max}2^{-n}$.

\emph{Step 2: Lipschitz dependence.} If two chords have endpoints at distance at most $\delta$ from the
corresponding endpoints of the other, then their lengths differ by at most $2\delta$: the length is
$|A-B|$, and $\bigl||A-B|-|A'-B'|\bigr|\le|A-A'|+|B-B'|$ by the triangle inequality.

\emph{Step 3: total.} Combining Steps 1 and 2 with $\delta=2R_{\max}2^{-n}$ gives an error of at most
$4R_{\max}2^{-n}$.
\end{proof}

\begin{remark}
With $R_{\max}\le5$ and $n=45$ the bound is about $5.7\times10^{-13}$, and with $n=60$ it is far below that. Since
the reported readings that decide the criterion are of order $10^{-9}$ and above, the evaluation error is four
orders of magnitude smaller than the smallest reported value. The readings are therefore not resolution
artefacts; the obstruction to turning them into a theorem is the covering radius, not the arithmetic.
\end{remark}

\subsection{Controls}

Two runs calibrate the method. On the unit disc, with $2000$ directions and the search domain $|P|\le0.5$, the
minimiser is located at the centre to within $|P^{*}|=0$ (all reported coordinates vanished) with
$V_{2000}\approx2.045\times10^{-27}$ and mean chord $2.0000000$; this is the positive control: when an exact zero
exists, the minimiser reaches machine zero, so the $10^{-9}$-level values seen on ellipses are not a floor imposed
by the method. The second control compares the fast chord evaluation with the closed-form chord of a circle,
$2\sqrt{1-d^2}$, and finds agreement to $6.35\times10^{-14}$; an earlier reference implementation using a
$2048$-edge polyline was found to degenerate near the origin and was discarded, which is why the later runs use
the analytic or bisection convention only.

\subsection{Ellipse family: readings}

Table~\ref{tab:ellipse-small} collects the first scan on five shapes: five ratios, $2000$ directions, an
$11\times11$ grid on $|P|\le0.6\min(a,b)$ followed by a compass search. Table~\ref{tab:ellipse-large} collects
the readings at the three large ratios, where the domain was enlarged and the interior was searched
independently. All values are values of $F$, not of $\sigma$; the threshold of
Definition~\ref{def:criterion} corresponds to $F<10^{-18}$.

\begin{table}[htbp]
\centering
\small
\caption{First scan on five shapes: $2000$ directions, $11\times11$ grid on $|P|\le0.6\min(a,b)$, compass search
to step $10^{-8}$. Column 3 is the smallest value of $F$ found; every value is strictly positive and far above
the threshold $10^{-18}$.}
\label{tab:ellipse-small}
\begin{tabular}{@{}cccc@{}}
\toprule
$a$ & $b$ & $\min F$ & verdict \\
\midrule
$2$ & $1$ & $3.0883\times10^{-2}$ & no candidate \\
$1.5$ & $1$ & $3.4659\times10^{-3}$ & no candidate \\
$3$ & $1$ & $2.10864\times10^{-1}$ & no candidate \\
$2.5$ & $1.5$ & $1.9917\times10^{-2}$ & no candidate \\
$1.1$ & $1$ & $1.0330\times10^{-5}$ & no candidate \\
\bottomrule
\end{tabular}
\end{table}

\begin{table}[htbp]
\centering
\small
\caption{Readings at the three large ratios. The domain ladder $0.9\to0.95\to0.99\to0.999$ never changed the
verdict; the argmin is interior and its distance from the centre is quoted in units of $\min(a,b)$; the last
column is the smallest value of $F$ found on that ratio, confirmed by a separate radial layer scan with $240$
angles per layer and by a fiftieth-digit recheck of the argmin.}
\label{tab:ellipse-large}
\begin{tabular}{@{}ccccc@{}}
\toprule
$a/b$ & domain & $|P^{*}|/\min(a,b)$ & interior? & $\min F$ \\
\midrule
$2.5$ & $|P|\le0.9$ & $0.92383$ & no (boundary) & $9.82997\times10^{-2}$ \\
$2.5$ & $|P|\le0.99$ & $0.92383$ & yes & $9.83\times10^{-2}$ \\
$3$ & $|P|\le0.99$ & $0.95174$ & yes & $2.11\times10^{-1}$ \\
$5$ & $|P|\le0.99$ & $0.98920$ & yes & $1.09$ \\
\bottomrule
\end{tabular}
\end{table}

Three further readings belong to this group. First, the nine-ratio scan with ratios
$\{1.01,1.05,1.1,1.3,1.5,2,2.5,3,5\}$, $b=1$, domain $|P|\le0.9\min(a,b)$, $2000$-direction search and
$8000$-direction recheck: the minimum over the nine ratios is $1.2254\times10^{-9}$, attained at $a/b=1.01$, and
every one of the nine minima is strictly positive. Second, the decisive recheck at $a/b=1.01$: the argmin is
interior, at $P^{*}=(-3.74\times10^{-10},-0.140371)$ with $|P^{*}|=0.140371$, and there
$V_{2000}=V_{8000}=1.22537006\times10^{-9}$, reproduced by the fiftieth-digit evaluation. Since
$\sigma=\sqrt{F}\approx3.5\times10^{-5}$, the covering-radius condition of
Proposition~\ref{prop:grid-certificate} fails by orders of magnitude for any grid used, so this reading excludes a
candidate on the grid but cannot exclude an equichordal point between grid points. Third, a log-least-squares fit
of the nine minima against $(a/b-1)$ gives $\min F\approx c\,(a/b-1)^{p}$ with $c=0.024915$, $p=3.514$ and
$R^2=0.9956$; this is a description of nine numbers and, as Section~\ref{sec:discussion} says, not an asymptotic law.

A separate run searched the whole interior of the ellipse in elliptic coordinates, with a compass search and extra
refinement near the boundary, in order to decide whether the value on the disc $|P|\le0.99\min(a,b)$ is the value
on the interior; the reading supports the disc value. It remains a reading on a grid. For seven of the located
minimisers a local recheck computed the Hessian and found it positive semidefinite; by
Lemma~\ref{lem:hessian} this gives a zero-free ball around each minimiser, which is a local statement and nothing
more.

\subsection{Odd-harmonic family: readings}

For the family of Definition~\ref{def:families} the origin is equichordal by Theorem~\ref{thm:odd-harmonic}, and
the recorded value $V_{2000}\approx6\times10^{-28}$ with constant chord length $2$ is consistent with the theorem,
as it must be. The question is whether a second point exists. A first attempt on the annulus
$|P|\in[0.01,0.9\,r_{\max}]$ with a $13\times13$ grid and eight starts collapsed back towards the origin and
answered nothing; the run is recorded as such. A hard annulus constraint, which forbids the local search from
returning to the origin disc, gives the readings of Table~\ref{tab:fourier}.

\begin{table}[htbp]
\centering
\small
\caption{Odd-harmonic family under a hard annulus constraint: $13\times13$ grid, eight compass starts, $2000$
directions, step floor $10^{-8}$; all infima are attained on the inner boundary of the annulus, and the shell
minimum is non-decreasing in $|P|$. The origin's variance, of order $10^{-28}$ to $10^{-32}$, serves as the
positive control.}
\label{tab:fourier}
\begin{tabular}{@{}cccc@{}}
\toprule
$\varepsilon$ & annulus & $\arg\min |P|$ & $\min F$ \\
\midrule
$10^{-3}$ & $[0.0100,\,0.9\,r_{\max}]$ & $0.0100$ & $6.9099\times10^{-10}$ \\
$3\times10^{-3}$ & $[0.0100,\,0.9\,r_{\max}]$ & $0.0100$ & $1.7186\times10^{-9}$ \\
$10^{-2}$ & $[0.0100,\,0.9\,r_{\max}]$ & $0.0100$ & $2.7846\times10^{-8}$ \\
$3\times10^{-2}$ & $[0.0100,\,0.9\,r_{\max}]$ & $0.0100$ & $2.9566\times10^{-7}$ \\
\bottomrule
\end{tabular}
\end{table}

A later run pushed the search inward: $13\times13$ grid, twelve radial rings, five compass starts, step floor
$10^{-9}$, over $|P|\in[10^{-6},0.9\,r_{\max}]$, with the $2000$- and $8000$-direction criteria at the argmin. No
second candidate was found, and the ring readings increase with $|P|$, from about $2.9\times10^{-16}$ at
$|P|=10^{-5}$ to about $1.5\times10^{-1}$ at $|P|=0.9$. The increase is the reason the constrained minima sit on the
inner boundary: the annulus is where the variance is smallest, and by Theorem~\ref{thm:odd-harmonic} the origin,
where it vanishes exactly, is just outside it.

A side reading calibrates the family itself: for coefficient sets whose radial function is not of odd-harmonic type,
the two-point residual $V(P_1)+V(P_2)$ under the constraint $|P_1-P_2|\ge0.05$ stays positive and grows with
$\varepsilon$ with slope $1.885$ and $R^2=0.9975$. This growth is a scaling of the residual and is not evidence
that a second equichordal point is absent; it is recorded because it shows that the odd-harmonic normalisation is
the special case, not the generic one.

\subsection{Runs that contributed nothing}

Three runs terminated abnormally or reported themselves inconclusive (two by a non-zero exit status, one by an
internal consistency check), and one run was a replication of another on the same grid. None of them contributes
evidence in either direction; the questions they were aimed at are answered, where they are answered at all, by the
runs listed above. A further run showed that the number of directions enters the evaluator non-trivially, which is
why both $V_{2000}$ and $V_{8000}$ are required in the criterion of Definition~\ref{def:criterion}.

\begin{proposition}[Quantified outcome of the finite scans]\label{prop:finite-scan}
Consider the finite set of data consisting of the five shapes of Table~\ref{tab:ellipse-small}, the nine ratios
of the nine-ratio scan, the three large ratios of Table~\ref{tab:ellipse-large}, and the four values of
$\varepsilon$ of Table~\ref{tab:fourier}, each evaluated on the grid, with the direction counts, domains and step
floors stated there. Then:
\begin{enumerate}
\item[(i)] for every grid point $P$ of every listed configuration, $F(P)>0$;
\item[(ii)] the smallest value attained anywhere in the listed configurations, namely $1.2254\times10^{-9}$ at
$a/b=1.01$, exceeds the threshold $10^{-18}$ by nine orders of magnitude;
\item[(iii)] consequently no point of any listed configuration satisfies the criterion of
Definition~\ref{def:criterion}, and the origin of the odd-harmonic family is the only point at which the recorded
variance is at the level of an exact zero.
\end{enumerate}
\end{proposition}

\begin{proof}
\emph{Step 1 (the quantifier is finite).} The proposition quantifies over a finite set: the grid points of finitely
many grids, evaluated at finitely many directions.

\emph{Step 2 (parts (i) and (ii) are readings of the tables).} Each of (i) and (ii) is a statement about the numbers
in Tables~\ref{tab:ellipse-small}, \ref{tab:ellipse-large} and~\ref{tab:fourier} together with the recorded minima
of the nine-ratio scan; the comparison $1.2254\times10^{-9}>10^{-18}$ is arithmetic.

\emph{Step 3 (part (iii) restates (ii)).} Part (iii) restates (ii) through Definition~\ref{def:criterion},
together with the observation that the variance computations at exact equichordal points return values of order
$10^{-27}$ or smaller, as the disc control and the origin of the odd-harmonic family both show.

\emph{Step 4 (scope).} Nothing in this proof is a statement about points that are not grid points, and by
Proposition~\ref{prop:grid-certificate} nothing of that kind could follow from these numbers.
\end{proof}

} % end of \CompEvidenceBody
\section{Discussion and open problems}\label{sec:discussion}

\subsection{What the two families suggest}

The scans are consistent with a clean picture, which we state as conjectures with their exact quantifiers and with
the list of what supports them.

\begin{conjecture}[Ellipses have no equichordal point]\label{conjecture-ellipse}
Let $a>b>0$. Then the ellipse $E(a,b)$ has no equichordal point in its interior.
\end{conjecture}

Support: the centre is excluded by Proposition~\ref{prop:ellipse-centre}; the scans of
Section~\ref{app:comp} cover ratios from $1.01$ to $5$, four domains from $0.6$ to $0.999$ times the small
semi-axis, one full-interior search, and two direction counts, and find no candidate. The conjecture is strictly
stronger than the quoted literature: the literature excludes two equichordal points on a restricted class of
curves, the conjecture excludes even one on an ellipse. Nothing in this manuscript proves it, and by
Proposition~\ref{prop:grid-certificate} the scans cannot, in their present form, prove it.

\begin{conjecture}[The odd-harmonic family has no second equichordal point]\label{conjecture-fourier}
Let $r(\theta)=1+\varepsilon(0.6111\cos3\theta+0.3889\cos5\theta)$ with $0<\varepsilon<1$. Then the origin is the only
equichordal point of the curve $\{r(\theta)(\cos\theta,\sin\theta)\}$.
\end{conjecture}

Support: Theorem~\ref{thm:odd-harmonic} establishes the origin; the annulus runs of
Table~\ref{tab:fourier} and the inward run down to $|P|=10^{-6}$ find no second candidate, and the ring readings
increase with $|P|$. The conjecture is stated for a single parametrised family with two harmonics and for a fixed
coefficient pair; it says nothing about other coefficient vectors of odd type, even though
Theorem~\ref{thm:odd-harmonic} shows that all of them have the origin as an equichordal point.

\subsection{The gap between the scans and the conjectures}

The following list says, for each conjecture, which step is missing. Items (1)--(4) are missing computations or
missing estimates of a known type; items (5)--(7) are missing ideas.

\begin{problem}[A global lower bound]\label{prob:global}
Find a positive lower bound for $F$ on a compact subdomain of $\operatorname{int}E(a,b)$ that does not rest on a
grid. By Proposition~\ref{prop:grid-certificate} it suffices to bound the modulus of continuity of $\sigma$ at the
covering radius used and compare; interval arithmetic or a branch-and-bound in the style of global optimisation
would give such a bound, and none has been carried out here.
\end{problem}

\begin{problem}[Ratios beyond five]\label{prob:ratios}
Extend the scans to $a/b>5$. The three large ratios scanned show the infimum increasing with $a/b$ and the argmin
approaching the boundary of the disc $|P|\le0.99\min(a,b)$, which is itself a warning: at large ratios the disc is
a small part of the interior, and a boundary argmin is exactly the case in which the reading does not localise the
minimum.
\end{problem}

\begin{problem}[The degenerate limit $a/b\to1$]\label{prob:limit}
Determine the behaviour of $\inf F$ as $a/b\to1$. The nine scanned minima decrease towards zero, and the fitted
power law $c\,(a/b-1)^{p}$ with $p\approx3.514$ and $R^2=0.9956$ describes the nine points; it is not an asymptotic
statement, and no uniform lower bound is available from it. What is needed is either a two-sided estimate of
$\inf F$ in terms of $(a/b-1)$ or a proof that the infimum is not attained in the limit.
\end{problem}

\begin{problem}[The limit $\varepsilon\to0$]\label{prob:eps}
For the odd-harmonic family only four values of $\varepsilon$ were scanned. Whether the absence of a second point
survives as $\varepsilon\to0$, and how the smallest annular variance scales in $\varepsilon$, is open here. The
family degenerates to the circle at $\varepsilon=0$, where every interior point at distance $d$ has
$F>0$ for $d>0$ by Proposition~\ref{prop:disc} and $F(0)=0$, so the $\varepsilon=0$ endpoint is understood and the
question is about uniformity in $\varepsilon$.
\end{problem}

\begin{problem}[The neighbourhood of the origin in the odd-harmonic family]\label{prob:origin}
The origin is equichordal, so $F$ vanishes there; a second point, if it existed, could in principle be arbitrarily
close to the origin. The constrained runs exclude an inner disc by construction at radius $0.01$, and the inward
run reaches $|P|=10^{-6}$ but only at isolated radii. A local analysis of $F$ near the origin, in the spirit of
Lemma~\ref{lem:hessian} but with the Hessian degenerate at an exact zero, is missing.
\end{problem}

\begin{problem}[Families that are not of odd-harmonic type]\label{prob:families}
Theorem~\ref{thm:odd-harmonic} uses only the relation $r(\theta+\pi)=2-r(\theta)$, which is one linear condition
on the Fourier coefficients. Whether a curve with an equichordal point necessarily satisfies a relation of this
kind, or whether equichordal curves with an equichordal point off the origin behave quite differently, is not
addressed. The quoted literature does provide symmetry properties of $2e$-curves independent of convexity, and
those are the natural starting point, but they are quoted at second hand and no proof was reconstructed here.
\end{problem}

\begin{problem}[The non-convex case]\label{prob:nonconvex}
The quoted negative result is stated for a class of curves that contains the convex curves, and the commentary
explicitly restricts it that way. Whether a non-star-shaped simple closed curve can carry two equichordal points
is not settled by anything quoted or proved here. Not finding a source that settles it is not the same as the
case being open, and we deliberately do not assert the latter.
\end{problem}

\subsection{On method}

Three methodological points are worth recording, because each of them cost a run to learn.

First, a positive control is not optional. The disc control shows that an exact equichordal point computes to
$10^{-27}$ or below; without that number, a reading of $10^{-9}$ on an ellipse could not be distinguished from a
resolution floor. Second, a negative control is not optional either: the discarded polyline reference would have
produced discrepancies near the origin that look exactly like the phenomenon being searched for, and the
comparison with the analytic circle chord is what exposed it. Third, a boundary argmin must be reported as such.
Four of the eleven shapes in the wider scan had their argmin on the boundary of the search disc; those readings
say that no zero was found inside the disc, not that the infimum inside the disc has been located, and they are
recorded with that qualification.

Finally, the distinction between what is proved and what is measured has to survive compression. Of the nine
statements with proofs in this manuscript, none asserts the non-existence of an equichordal point on a family of
curves; the two statements that do assert non-existence of an equichordal point (one for a disc, one for the
centre of an ellipse) are proved directly. The two conjectures assert non-existence over families, and they are
supported by finite scans with an explicit covering-radius obstruction. No proof in this manuscript has been
machine-checked, and no statement here claims to have been.

\begin{thebibliography}{9}

\bibitem{klee1960}
V.\ Klee,
\emph{A problem collection}, in: Convexity (Proceedings of Symposia in Pure Mathematics), 1960; reprinted
edition consulted. The problem statement on printed p.\,24 and the commentary on printed p.\,27 are the sources of
the verbatim quotations in Section~\ref{sec:intro-lit}.

\bibitem{fujiwara1916}
T.\ Fujiwara, 1916. The attribution of the question to Fujiwara is quoted here from the commentary in
\cite{klee1960}; the original article was not consulted.

\bibitem{blaschke1917}
W.\ Blaschke, E.\ Rothe, R.\ Weitzenb\"ock, 1917. The independent attribution of the question is quoted here from
the commentary in \cite{klee1960}; the original articles were not consulted.

\bibitem{rychlik1996}
J.\ Rychlik, \emph{The equichordal point problem}, 1996. Registered as a source for this work but not available in
full text, so no statement of this manuscript rests on it.

\bibitem{rychlik1997}
J.\ Rychlik, \emph{A complete solution to the equichordal point problem of Fujiwara, Blaschke, Rothe and
Weitzenb\"ock}, Inventiones Mathematicae \textbf{129} (1997), no.\ 1, 141--212. The result is used here only at the
strength of the second-hand report quoted in Section~\ref{sec:intro-lit}, namely for a class of curves that
includes the convex curves; the paper itself was not available in a machine-readable text and nothing in it is
paraphrased.

\end{thebibliography}

\appendix

\section{Computational evidence}\label{app:comp}

\CompEvidenceBody

\subsection{Reproducibility index}\label{app:index}

The following table maps each statement of this manuscript to the records that support it. It is the only place in
this manuscript where internal identifiers appear.
\begin{center}
\small
\begin{tabular}{@{}p{4.2cm}p{9.6cm}@{}}
\toprule
object & record \\
\midrule
Definitions of equichordal point and equichordal curve & stmt:000000, stmt:000001; source lit:000007 (verbatim, quote\_verified=true) \\
Scope convention: convex body versus simple closed curve & stmt:000002; source lit:000000 at bibliographic level only \\
Verbatim quotations of Section~\ref{sec:intro-lit} & rel:000000 (printed p.\,24), rel:000001 (printed p.\,24), rel:000002, rel:000003, rel:000004, rel:000005 (printed p.\,27); source file library/000007.txt \\
Theorem~\ref{thm:odd-harmonic} & natural-language proof, no experiment needed; numerical consistency check exp:000002 (variance $\approx6\times10^{-28}$, chord $2$) \\
Proposition~\ref{prop:disc} & exp:000001 (centre located, $V_{2000}\approx2.045\times10^{-27}$); method check exp:000003 (bisection versus analytic chord, $6.35\times10^{-14}$) \\
Proposition~\ref{prop:ellipse-centre} & proof only; consistent with exp:000000 \\
Table~\ref{tab:ellipse-small} & exp:000000 \\
Table~\ref{tab:ellipse-large} & exp:000011 (domain ladder), exp:000012 (interior minima), exp:000015 (full interior), exp:000017 (radial layers) \\
Nine-ratio scan and power-law fit & exp:000004 (nine ratios, minimum $1.2254\times10^{-9}$, $c=0.024915$, $p=3.514$, $R^2=0.9956$); wider scan exp:000008 (eleven shapes, four boundary argmins) \\
Decisive recheck at $a/b=1.01$ & exp:000010 \\
Local Hessian rechecks & exp:000018 (seven minimisers) \\
Table~\ref{tab:fourier} & exp:000006 (hard annulus), exp:000007 (decisive annulus run) \\
Inward run to $|P|=10^{-6}$ & exp:000019 \\
Runs contributing nothing & exp:000005 (collapsed towards the origin), exp:000009 and exp:000016 (non-zero exit status), exp:000013 (self-reported inconclusive), exp:000014 (reading of exp:000013); exp:000020 is a replication of exp:000018 and is counted once \\
Conjecture~\ref{conjecture-ellipse} & conjectures 000000, status supported\_not\_proved \\
Conjecture~\ref{conjecture-fourier} & conjectures 000001, status supported\_not\_proved \\
\bottomrule
\end{tabular}
\end{center}

All numerical records carry evidence level cross\_check in the sense of this workflow: independent
re-computation, not proof. Scripts, parameters and the recorded content digests of the scripts are stored with the
experiment records in the project directory; the aggregate numerical output of the ellipse scans is in
out/eq\_s1/results.json and the run logs are in out/experiments/.

\end{document}
